\documentclass[11pt,a4paper]{article}
\usepackage[T1]{fontenc}
\usepackage[utf8]{inputenc}
\usepackage{lmodern}
\usepackage{amsmath,amssymb,amsthm}
\usepackage{booktabs}
\usepackage{xcolor}
\usepackage[margin=27mm]{geometry}
\usepackage[colorlinks=true,urlcolor=blue,citecolor=blue,linkcolor=blue]{hyperref}
\newtheorem{theorem}{Theorem}
\newtheorem{proposition}[theorem]{Proposition}
\newtheorem{lemma}[theorem]{Lemma}
\newtheorem{corollary}[theorem]{Corollary}
\theoremstyle{definition}
\newtheorem{definition}[theorem]{Definition}
\newtheorem{remark}[theorem]{Remark}
\theoremstyle{plain}
\newcommand{\ket}[1]{\lvert#1\rangle}
\newcommand{\bra}[1]{\langle#1\rvert}
\newcommand{\norm}[1]{\lVert#1\rVert}
\newcommand{\one}{\mathbf{1}}
\newcommand{\R}{\mathbb{R}}
\newcommand{\E}{\mathbb{E}}
\newcommand{\T}{\mathsf{T}}
\newcommand{\ip}[2]{\langle#1,#2\rangle}
\setlength{\emergencystretch}{2em}
\setlength{\parindent}{0pt}
\setlength{\parskip}{0.3em plus 0.1em minus 0.1em}
\title{The variance of a sum of independent quantum errors:\\
a proof for centred errors and counterexamples without centring}
\author{Fable \& Astra\\[0.35em]
\normalsize Research led by Jean-Philippe Bruneton\\[-0.1em]
\small Universit\'e Paris Cit\'e, France}
\date{19 November 2026}

\begin{document}
\maketitle
\begin{abstract}
When a quantum computation suffers several independent errors, one would like
to predict the size of the total error from the sizes of its parts. Lacalle
and Pozo Coronado model an error as a random pure state scattered around the
error-free state on the unit sphere of the state space, measure its size by
the mean squared distance $V$ to the error-free state, and prove that two
independent isotropic errors combine as
$V_{12}=V_1+V_2-\tfrac12V_1V_2$. They conjecture the same formula for
arbitrary independent errors and state it as a theorem on the circle. We show
that the answer depends on one hypothesis that the conjecture leaves implicit:
whether the errors are centred, in the sense that their Euclidean mean points
along the error-free state. The formula is equivalent to the multiplicativity
of the mean overlap $m=1-V/2$, and composing independent errors multiplies
mean operators. Hence the formula holds for any number of independent centred
errors, isotropic or not, in both natural definitions of the sum, and an
exact identity gives the correction term in general. The multiplicativity
itself is elementary and classical; the contribution of the note is to
identify centring as the hypothesis on which the conjecture turns and to
supply exact counterexamples. Without centring the
formula fails: two independent errors on the circle with density
$(1+\sin\theta)/(2\pi)$, which belongs to the family used in the source's own
proof, give $V_{12}=5/2$ instead of $2$, so the theorem on the circle is false
as stated; two deterministic over-rotations $R_x(2\pi/3)$ of one qubit give
$V_{12}=3$ instead of $3/2$. The same mechanism proves the companion fidelity
formula of Lacalle, Pozo Coronado and Mart\'in-Cuevas when the second error
has an isotropic second moment.
\end{abstract}

\section{Introduction}
Errors in a quantum computer accumulate. A gate that misses its target by a
small amount leaves the register in a pure state close to, but not equal to,
the intended one, and a long computation composes many such imperfections.
A simple model treats each error as a random point on the unit sphere of
state space, concentrated around the error-free state, and asks how the
spread of the final state depends on the spread of the individual errors.
Because the sphere is compact, spreads cannot simply add: a total error never
exceeds the diameter of the sphere. Lacalle and Pozo
Coronado~\cite{LacallePozo2019} measured the spread by the mean squared
distance to the error-free state, which they call the variance, and found that
for two independent errors whose densities depend only on the distance to
their centre (isotropic errors) the variances combine by the rule
\begin{equation}\label{eq:formula}
V(X_1+X_2)=V(X_1)+V(X_2)-\frac{V(X_1)\,V(X_2)}{2}.
\end{equation}
They conjectured that the rule holds for any two independent errors, proved it
for arbitrary independent random variables on the circle, and illustrated the
general case by an exact computation on the seven-dimensional sphere of two
qubits.

This note settles the conjecture by identifying the hypothesis on which it
turns. The variance of an error depends only on the Euclidean mean of the
random state, and composing independent errors multiplies mean operators.
Rule~\eqref{eq:formula} is the statement that the mean overlap with the
error-free state is multiplicative, and this holds whenever the error-free
state is an eigenvector of the mean operators, that is, whenever the errors
are \emph{centred}: their Euclidean mean lies along the error-free state.
Isotropy guarantees centring but is far stronger than needed. The source paper
itself assumes that the error-free state ``represents the mean'' of each
error, and under that convention the conjecture is true for anisotropic errors
with different distributions, any number of errors and every number of
qubits, with a three-line proof (Theorems~\ref{thm:gate}
and~\ref{thm:transport}, Corollary~\ref{cor:k}). Without centring the rule is
false, and the theorem on the circle fails already for a pair of densities
from the family used in its proof (Proposition~\ref{prop:counter}). The
physically important case is that of biased coherent errors, such as a
systematic over-rotation, for which the variance grows faster than the rule
predicts.

The simplest illustration is a single qubit subject twice to the same
deterministic over-rotation $R_x(2\pi/3)=\exp(-i\tfrac\pi3\sigma_x)$. Each
rotation has overlap $\tfrac12$ with $\ket0$, hence variance $1$, and
rule~\eqref{eq:formula} predicts $\tfrac32$ for the pair. The composite
rotation $R_x(4\pi/3)$ has overlap $-\tfrac12$, hence variance $3$. The
discrepancy comes from the transverse component of the mean, which is
nonzero for a biased rotation and absent for a centred error; the general
identity~\eqref{eq:identity} below measures exactly this component.

\section{Related work}\label{sec:related}
The model and the conjecture are those of Lacalle and Pozo
Coronado~\cite{LacallePozo2019}, published in 2019 and posted to arXiv in
December 2024. They prove rule~\eqref{eq:formula} for independent isotropic
errors on every sphere $S^d$ (their Theorem~2), through an explicit expansion
of the densities in powers of $\cos\theta_0$, double-factorial integrals and a
combinatorial identity on ordered trees. Their Theorem~3 asserts the rule for
any two independent random variables on $S^1$. Its proof expands densities in
the functions $g_k(\theta)\propto1+\cos^k\theta$ and
$h_k(\theta)\propto1+\sin^k\theta$, treats the pairs $(g_k,g_l)$, $(g_k,h_l)$
and $(h_k,g_l)$, and omits the pair $(h_k,h_l)$; Proposition~\ref{prop:counter}(a)
shows that this pair violates the rule. Their Example on $S^7$ combines an
anisotropic, non-centred first error with an isotropic second error and
obtains $V=445/224$; Section~\ref{sec:example} explains why this example
falls under Theorem~\ref{thm:transport} and therefore does not test the
anisotropic second error.

Lacalle, Pozo Coronado and Mart\'in-Cuevas~\cite{LacallePozoMartin2026}
recall in 2026 that rule~\eqref{eq:formula} ``was proven for isotropic
quantum errors and conjectured for general quantum errors'' and pose a twin
conjecture for the fidelity. Remark~\ref{rem:fidelity} proves that twin
formula in the transport model when the second error has an isotropic second
moment. We found no proof, refutation or correction of the variance
conjecture in the literature examined (Section~\ref{sec:scope}).

The distinction between centred and biased errors parallels the familiar
contrast between stochastic and coherent noise in fault tolerance: errors of
equal average size can have very different worst-case effects, and coherent
errors appear to produce the largest gap between average-case and worst-case
error~\cite{Kueng2016}. Under composition, moreover, the infidelity of
purely stochastic errors grows linearly in the number of gates, whereas that
of purely unitary errors can grow quadratically, and the average gate
fidelity of a composite channel is bounded between these two
regimes~\cite{CarignanDugas2019}. Proposition~\ref{prop:counter}(b) is an
instance of this coherent accumulation in the variance language
of~\cite{LacallePozo2019}.

The mechanism behind our positive results is classical in substance. On the
circle, independence of two angles gives
$\E e^{i(\theta_1+\theta_2)}=\E e^{i\theta_1}\,\E e^{i\theta_2}$, the
multiplicativity of the first trigonometric moment of a sum of independent
angles, whose real part is the identity used in
Proposition~\ref{prop:counter}(a). In general, the mean operators of
independent random isometries multiply,
$\E[R_2R_1]=\E[R_2]\,\E[R_1]$, which is the defining-representation case of
the fact that convolution of measures on a compact group multiplies their
Fourier transforms. We claim no novelty for these facts; what the note adds
is their application to the conjecture of~\cite{LacallePozo2019}, the
identification of centring as the relevant hypothesis, and the exact
correction term~\eqref{eq:identity}.

\section{Setting}\label{sec:setting}
A pure state of $n$ qubits, $\Psi=(x_0+ix_1,\dots,x_{d-1}+ix_d)$ in the
computational basis, is a point of the real unit sphere
$S^d\subset\R^{d+1}$ with $d=2^{n+1}-1$. We write $\ip\cdot\cdot$ for the
Euclidean inner product of $\R^{d+1}$; for two states,
$\ip xy=\operatorname{Re}\langle\Psi_x|\Psi_y\rangle$. The error-free state is
$P=(1,0,\dots,0)$, that is $\ket0$. An error is a random variable $X$ on
$S^d$, and its variance in the sense of~\cite{LacallePozo2019} is
\[
V(X)=\E\norm{X-P}^2=2-2\,\E[x_0]\in[0,4].
\]
We define the \emph{mean overlap} $m(X)=\E\ip PX=\E[x_0]=1-V(X)/2$. Since
$1-\tfrac12\bigl(V_1+V_2-\tfrac12V_1V_2\bigr)=\bigl(1-\tfrac12V_1\bigr)
\bigl(1-\tfrac12V_2\bigr)$, rule~\eqref{eq:formula} is equivalent to
\begin{equation}\label{eq:mult}
m(X_1+X_2)=m(X_1)\,m(X_2).
\end{equation}

\begin{definition}
An error $X$ is \emph{centred on $P$} if its Euclidean mean is parallel to
$P$: $\E[X]=m(X)\,P$, that is, $\E[x_j]=0$ for $1\le j\le d$.
\end{definition}

This is the precise content of the convention in~\cite{LacallePozo2019} that
$P$ ``represents the mean value'' of the distribution. An isotropic error is
centred: its law is invariant under the isometries of $\R^{d+1}$ that fix $P$
(the rotations about the axis $P$ when $d\ge2$, the reflection
$\theta\mapsto-\theta$ when $d=1$), so $\E[X]$ is fixed by them and lies in
$\R P$. The converse fails, since centring imposes $d$ linear conditions on
the law, whereas isotropy prescribes its whole angular dependence.

The source defines the sum $X_1+X_2$ explicitly only when the second error is
isotropic: the second error is then applied around the point reached by the
first, with a density depending only on the distance to that point. Two
definitions extend this to anisotropic errors.

\emph{Gate model.} Each error is a random isometry $R_j\in O(d+1)$; for $n$
qubits, a random unitary in $U(2^n)\subset O(2^{n+1})$, for instance a noisy
gate composed with the inverse of the ideal gate. The isometries $R_1,R_2$ are
independent, $X_j=R_jP$ is the error seen from the reference state, and
$X_1+X_2:=R_2R_1P$. This is the physical definition: the second error acts as
an operator on the already erroneous state.

\emph{Transport model.} The first error is a random point $X_1\in S^d$; the
second is a random point $Y_2\in S^d$, independent of $X_1$, whose law is
given around $P$. Choose a measurable section $T:S^d\to O(d+1)$ with
$T(x)P=x$ and set $X_1+X_2:=T(X_1)\,Y_2$: the law of the second error is
copied around the point reached by the first. When $Y_2$ is isotropic the law
of $T(x)Y_2$ does not depend on the choice of $T$, and this is the definition
of~\cite{LacallePozo2019}.

\section{Results}\label{sec:results}
\subsection{The identity behind the formula}

\begin{theorem}[gate model]\label{thm:gate}
Let $R_1,R_2$ be independent random isometries in $O(d+1)$, $X_j=R_jP$,
$m_j=m(X_j)$ and $V_j=V(X_j)$. Let $u_1=\E[X_1]-m_1P$ be the transverse mean
of the first error and $w_2=\E[R_2^\T P]-m_2P$ the transverse mean of the
inverse error $R_2^{-1}P$. Then
\begin{equation}\label{eq:identity}
V(X_1+X_2)=V_1+V_2-\frac{V_1V_2}{2}-2\,\ip{w_2}{u_1}.
\end{equation}
Rule~\eqref{eq:formula} therefore holds if and only if $\ip{w_2}{u_1}=0$; in
particular it holds for every second error as soon as the first error is
centred.
\end{theorem}

\begin{proof}
We have $\ip P{R_2R_1P}=\ip{R_2^\T P}{R_1P}$. The random vectors $R_2^\T P$
and $R_1P$ are bounded and independent, and the inner product is bilinear, so
$m(X_1+X_2)=\ip{\E[R_2^\T P]}{\E[R_1P]}$. Moreover
$\ip P{\E[R_2^\T P]}=\E\ip{R_2P}P=m_2$, so $\E[R_2^\T P]=m_2P+w_2$ with
$w_2\perp P$, and likewise $\E[R_1P]=m_1P+u_1$ with $u_1\perp P$. Hence
$m(X_1+X_2)=m_1m_2+\ip{w_2}{u_1}$, and $V=2-2m$ gives~\eqref{eq:identity}.
\end{proof}

\begin{theorem}[transport model]\label{thm:transport}
Let $X_1$ be any random point of $S^d$, let $Y_2$ be a random point of $S^d$
independent of $X_1$ and centred on $P$, and let $T$ be any measurable
section. Then $\E[T(X_1)Y_2\mid X_1]=m(Y_2)\,X_1$, and rule~\eqref{eq:formula}
holds for $X_1+X_2=T(X_1)Y_2$.
\end{theorem}

\begin{proof}
Conditionally on $X_1$ the matrix $T(X_1)$ is fixed and $Y_2$ keeps its law,
so $\E[T(X_1)Y_2\mid X_1]=T(X_1)\E[Y_2]=m(Y_2)\,T(X_1)P=m(Y_2)\,X_1$. Taking
expectations and the inner product with $P$ gives
$m(X_1+X_2)=m(Y_2)\,m(X_1)$, which is~\eqref{eq:mult}.
\end{proof}

\begin{corollary}[the conjecture under the centring convention]\label{cor:k}
Let $X_1,\dots,X_k$ be mutually independent errors, each centred on $P$,
isotropic or not, with equal or different laws. In either model,
\[
1-\frac{V(X_1+\dots+X_k)}{2}=\prod_{j=1}^{k}\Bigl(1-\frac{V(X_j)}{2}\Bigr).
\]
In particular rule~\eqref{eq:formula} holds for $k=2$ in every dimension
$d\ge1$, hence for every number of qubits.
\end{corollary}

\begin{proof}
In the gate model, centring means $\E[R_j]P=m_jP$. By mutual independence
$\E[R_k\cdots R_1]=\E[R_k]\cdots\E[R_1]$, and
$\E[R_k]\cdots\E[R_1]P=m_k\cdots m_1P$ by induction; take the inner product
with $P$. In the transport model, let $X^{(j)}$ be the state after $j$
errors. Theorem~\ref{thm:transport}, applied conditionally on $X^{(j-1)}$,
gives $\E[X^{(j)}\mid X^{(j-1)}]=m_jX^{(j-1)}$, hence
$\E[X^{(k)}]=m_k\cdots m_1P$ by the tower property.
\end{proof}

The main theorem of~\cite{LacallePozo2019} is the special case of
Corollary~\ref{cor:k} in which every law is invariant under the isometries
fixing $P$. Isotropy is used there only to make $P$ an eigenvector of the
mean operators, and centring is enough for that.

\subsection{Counterexamples without centring}

\begin{proposition}\label{prop:counter}
Rule~\eqref{eq:formula} fails for independent errors that are not centred, in
each of the following exact situations.
\begin{enumerate}
\item[(a)] \emph{On the circle.} Let $X_1,X_2$ be independent with density
$h_1(\theta)=(1+\sin\theta)/(2\pi)$ on $S^1$. Then $V(X_j)=2$, so
rule~\eqref{eq:formula} predicts $2$, whereas $V(X_1+X_2)=5/2$ when the sum
is the rotation $\theta=\theta_1+\theta_2$, and $3/2$ with the convention
$\theta=\theta_2-\theta_1$. In either case Theorem~3
of~\cite{LacallePozo2019} is false as stated.
\item[(b)] \emph{One qubit, coherent errors.} For $U_1=U_2=R_x(2\pi/3)$ one has
$V_1=V_2=1$ and $V(X_1+X_2)=3$, whereas rule~\eqref{eq:formula} predicts
$3/2$.
\item[(c)] \emph{Centring the second error is not enough in the gate model.}
On one qubit, let $U_2$ equal $R_y(\pi/2)$ or $R_y(-\pi/2)Z$ with probability
$\tfrac12$ each, and let $X_1=R_y(\pi/2)\ket0$ be deterministic. Then the
second error is centred with $m_2=\tfrac{\sqrt2}2$, $m_1=\tfrac{\sqrt2}2$, and
$m(X_1+X_2)=0\ne\tfrac12$, so $V(X_1+X_2)=2$ instead of $1$.
\end{enumerate}
\end{proposition}

\begin{proof}
(a) The density $h_1$ is nonnegative with total mass $1$, and
$\E\cos\theta=0$, $\E\sin\theta=\tfrac12$. Hence $m(X_j)=0$ and $V(X_j)=2$.
By independence,
$\E\cos(\theta_1\pm\theta_2)=\E\cos\theta_1\E\cos\theta_2\mp
\E\sin\theta_1\E\sin\theta_2=\mp\tfrac14$, so $V(X_1+X_2)=2\pm\tfrac12$.
The density $h_1$ belongs to the family $\{h_k\}$ used in the proof of
Theorem~3 of~\cite{LacallePozo2019}; the case $(h_k,h_l)$ is the one missing
from that proof. On $S^1$ the correct statement is
$V(X_1+X_2)=V_1+V_2-\tfrac12V_1V_2+2\,\E[\sin\theta_1]\,\E[\sin\theta_2]$ for
the rotation convention, the special case of~\eqref{eq:identity} with
$u_1=\E[\sin\theta_1]e_1$ and $w_2=-\E[\sin\theta_2]e_1$.

(b) In the real picture, $\bra0R_x(\alpha)\ket0=\cos(\alpha/2)$, so
$m_1=m_2=\cos(\pi/3)=\tfrac12$ and $V_1=V_2=1$. The product is
$R_x(4\pi/3)$, with $m(X_1+X_2)=\cos(2\pi/3)=-\tfrac12$ and
$V(X_1+X_2)=3$. More generally, for independent random over-rotations
$R_x(\alpha_j)$ the identity~\eqref{eq:identity} reads
$V_{12}=V_1+V_2-\tfrac12V_1V_2+2\,\E[\sin\tfrac{\alpha_1}2]\,
\E[\sin\tfrac{\alpha_2}2]$: a systematic bias of the same sign makes the
variance grow faster than rule~\eqref{eq:formula}, and the rule is exact when
one of the angles has a symmetric law.

(c) With $R_y(\alpha)=\begin{pmatrix}\cos\frac\alpha2&-\sin\frac\alpha2\\
\sin\frac\alpha2&\cos\frac\alpha2\end{pmatrix}$ and $Z=\mathrm{diag}(1,-1)$,
the mean operator is
$M_2=\E[U_2]=\tfrac{\sqrt2}{2}\begin{pmatrix}1&-1\\0&0\end{pmatrix}$, so
$M_2\ket0=\tfrac{\sqrt2}2\ket0$ and the second error is centred, whereas
$M_2^\T\ket0=\tfrac{\sqrt2}2(\ket0-\ket1)$ has a transverse part. With
$X_1=\tfrac{\sqrt2}2(\ket0+\ket1)$, independence gives
$m(X_1+X_2)=\bra0M_2X_1=\tfrac12(1-1)=0$, while
$m_1m_2=\tfrac12$. This agrees with~\eqref{eq:identity}:
$u_1=\tfrac{\sqrt2}2\ket1$, $w_2=-\tfrac{\sqrt2}2\ket1$ and
$\ip{w_2}{u_1}=-\tfrac12$.
\end{proof}

In the gate model it is therefore the centring of the first error, or of the
inverse second error $R_2^{-1}P$, that matters. In the transport model, when
the second error is not centred, the variance of the sum depends in addition
on the choice of the section $T$, so that the sum is not determined by the
laws of $X_1$ and $X_2$ alone; the checker of Section~\ref{sec:verification}
exhibits finite ensembles for which two sections give variances differing by
more than $1$.

\subsection{The example on the seven-sphere}\label{sec:example}
The example illustrating the conjecture in~\cite{LacallePozo2019} uses
hyperspherical angles $(\theta_0,\theta_1,\theta_2,\dots)$ on $S^7$, with
$x_0=\cos\theta_0$ and $x_1=\sin\theta_0\cos\theta_1$, and takes the
densities
\[
g_1\propto2+\cos\theta_0\sin^2\theta_1+\cos\theta_1\sin\theta_2,\qquad
g_2\propto1+\cos\theta_0
\]
for the first and the second error. The density $g_2$ is isotropic, hence
centred. The density $g_1$ is anisotropic and not centred: an exact
integration gives $\E[x_1]=32768/(55125\,\pi^2)\approx0.0602$. The example
therefore falls under Theorem~\ref{thm:transport}, which gives
$V=2-2m_1m_2$ with $m_1=\tfrac3{56}$ and $m_2=\tfrac18$, that is
$V=\tfrac{445}{224}\approx1.98661$, the value obtained
in~\cite{LacallePozo2019} by a sixfold integral. It tests neither an
anisotropic nor a non-centred second error.

\subsection{The companion fidelity formula}

\begin{remark}\label{rem:fidelity}
Lacalle, Pozo Coronado and Mart\'in-Cuevas~\cite{LacallePozoMartin2026}
conjecture, for $F^2(X)=\E[x_0^2+x_1^2]=\E\lvert\langle0|\Psi_X\rangle\rvert^2$
and $f_0^2=2/(d+1)$, that
\[
F^2(X_1+X_2)=f_0^2+\frac{\bigl(F^2(X_1)-f_0^2\bigr)\bigl(F^2(X_2)-f_0^2\bigr)}{1-f_0^2}.
\]
In the transport model this holds, for an arbitrary first error and any
section, whenever the second error has an isotropic second moment,
$\E[Y_2Y_2^\T]=a\,PP^\T+b\,(\one-PP^\T)$ with $a+db=1$.
\end{remark}

\begin{proof}
Conditionally on $X_1$,
$\E[XX^\T\mid X_1]=T(X_1)\E[Y_2Y_2^\T]T(X_1)^\T
=a\,X_1X_1^\T+b\,(\one-X_1X_1^\T)$, independently of the section. Taking the
sum of the $(0,0)$ and $(1,1)$ entries and then expectations,
$F^2(X_1+X_2)=(a-b)F^2(X_1)+2b$. Moreover $F^2(X_2)=a+b$, and $a+db=1$
implies $1-(a-b)=(d+1)b$, so that $2b=f_0^2\bigl(1-(a-b)\bigr)$ and
$a+b-f_0^2=(a-b)(1-f_0^2)$. Substituting gives the formula.
\end{proof}

The checker shows that the fidelity formula fails for a generic second error
in the transport model. The gate model is not treated here.

\section{Verification}\label{sec:verification}
The statements above are proved analytically. The public
\href{https://github.com/jpbruneton/open_quantum_problems/blob/main/research/entries/variance-sum-quantum-errors/checks/check_variance_formula.py}{Python checker}
tests them, and its
\href{https://github.com/jpbruneton/open_quantum_problems/blob/main/research/entries/variance-sum-quantum-errors/data/verification.json}{recorded output}
is part of the same entry. In exact arithmetic (SymPy) it recomputes the three
counterexamples of Proposition~\ref{prop:counter} and the example of
Section~\ref{sec:example}: $V(g_1)=\tfrac{53}{28}$, $V(g_2)=\tfrac74$,
$2-2m_1m_2=\tfrac{445}{224}$ and $\E[x_1]\ne0$ under $g_1$. A Monte Carlo
simulation of that example in the transport model, with about $3\times10^6$
samples and a fixed seed, agrees with $\tfrac{445}{224}$ within four standard
errors and excludes the value $2$. In floating-point arithmetic the checker
verifies the identity~\eqref{eq:identity} to $10^{-12}$ on $360$ random finite
ensembles of anisotropic, non-identically distributed isometries in $O(2)$,
$O(4)$ and $O(8)$, the exactness of rule~\eqref{eq:formula} when the first
error is centred and its generic failure (deviations up to $0.67$),
Theorem~\ref{thm:transport} on $45$ finite ensembles on $S^1$, $S^3$ and $S^7$
with two different sections, and Remark~\ref{rem:fidelity} on $S^3$ and $S^7$.
The run takes a few seconds.

\section{Scope and limitations}\label{sec:scope}
\emph{Analytic proof.} Theorems~\ref{thm:gate} and~\ref{thm:transport},
Corollary~\ref{cor:k}, Proposition~\ref{prop:counter} and
Remark~\ref{rem:fidelity} are proved in full; the remark only in the
transport model.

\emph{What this says about the conjecture.} The conjecture specifies neither
the sum of anisotropic errors nor whether centring is assumed. Read with the
centring convention stated in~\cite{LacallePozo2019}, it is true in both
natural models, and the proof uses neither isotropy nor equality of the laws.
Read without that convention, as the theorem on the circle and the example on
$S^7$ suggest, since both use non-centred densities, it is false, and the
theorem on the circle must be restricted to pairs in which one error is
centred. The weakest sufficient condition in the gate model is
$\ip{w_2}{u_1}=0$. Independence enters only through the factorisation
$\E[R_2^\T P\otimes R_1P]=\E[R_2^\T P]\otimes\E[R_1P]$; for $k$ errors the
expectation of the full product must factorise, which pairwise independence
does not ensure.

\emph{Physical scope.} Centring excludes biased coherent errors, such as a
systematic over-rotation, for which the variance grows faster than
rule~\eqref{eq:formula} predicts, as in
Proposition~\ref{prop:counter}(b). The threshold estimate that the source
derives from the rule, a variance per gate below $C/k$ for $k$ gates, relies on
the absence of such bias. The variance $V$ is sensitive to the global phase;
the ``quantum variance'' $V_q=2-2\,\E\sqrt{x_0^2+x_1^2}$ introduced
in~\cite{LacallePozo2019} to remove this dependence is not a function of
$\E[X]$, and none of the results above concerns it. The two models of the sum
cover the natural readings of the source without exhausting every
conceivable definition, for instance a Markov kernel whose shape depends on
the starting point.

\emph{Bibliographic priority.} The source was read in its arXiv \TeX{} file,
and the 2026 companion paper in its published version. The works citing the
source were listed through the Semantic Scholar interface and complemented by
web searches and an OpenAlex query; none of them proves, refutes or corrects
the conjecture, and its authors still describe it as conjectural in January
2026. Priority is claimed only for the application to this conjecture, the
role of centring, the identity~\eqref{eq:identity} and the counterexamples,
within the limits of this documented search. The underlying multiplicativity
of trigonometric moments and of mean operators is classical
(Section~\ref{sec:related}), and the isotropic case of the source is a
special case of that multiplicativity.

\begingroup\small
\begin{thebibliography}{9}
\setlength{\itemsep}{2pt}
\bibitem{LacallePozo2019}
J. Lacalle and L. M. Pozo Coronado,
\emph{Variance of the sum of independent quantum computing errors},
Quantum Inf. Comput. \textbf{19}, 1294--1312 (2019);
\href{https://doi.org/10.26421/QIC19.15-16-3}{doi:10.26421/QIC19.15-16-3};
\href{https://arxiv.org/abs/2412.15800}{arXiv:2412.15800}.

\bibitem{LacallePozoMartin2026}
J. Lacalle, L. M. Pozo Coronado and R. Mart\'in-Cuevas,
\emph{Fidelity of the sum of independent quantum computing errors},
Quantum Inf. Process. \textbf{25}(2), 38 (2026);
\href{https://doi.org/10.1007/s11128-025-05037-5}{doi:10.1007/s11128-025-05037-5}.

\bibitem{Kueng2016}
R. Kueng, D. M. Long, A. C. Doherty and S. T. Flammia,
\emph{Comparing Experiments to the Fault-Tolerance Threshold},
Phys. Rev. Lett. \textbf{117}, 170502 (2016);
\href{https://doi.org/10.1103/PhysRevLett.117.170502}{doi:10.1103/PhysRevLett.117.170502};
\href{https://arxiv.org/abs/1510.05653}{arXiv:1510.05653}.

\bibitem{CarignanDugas2019}
A. Carignan-Dugas, J. J. Wallman and J. Emerson,
\emph{Bounding the average gate fidelity of composite channels using the unitarity},
New J. Phys. \textbf{21}, 053016 (2019);
\href{https://doi.org/10.1088/1367-2630/ab1800}{doi:10.1088/1367-2630/ab1800};
\href{https://arxiv.org/abs/1610.05296}{arXiv:1610.05296}.
\end{thebibliography}
\endgroup
\end{document}
